The ball identity follows immediately from the two definitions:
\[
 \mathfrak B_{f,e}(P)
 =\{Q:P\ll Q,\ Q\ll P,\ D_f(P\|Q)\le B_f(e)\}
 =\mathfrak B^{\to}_{f,e}(P)\cap\{Q:Q\ll P\}. \tag{14a}
\]
For the mixture classes, if \(Q=eP+(1-e)R\), then \((1-e)R\le Q\).  Since \(1-e>0\),
\[
 Q\ll P\quad\Longleftrightarrow\quad R\ll P,
\]
where the reverse direction also uses \(eP\ll P\).  Consequently
\[
 \mathfrak L_e(P)=\mathfrak L^{\to}_e(P)\cap\{Q:Q\ll P\}. \tag{14b}
\]
Finally, consistency gives \(P=\mathcal L(Y(a)\mid A=a)\), and the one-sided completeness theorem gives \(P\ll Q\) for every \(Q\in\mathfrak M^{\to}_e(P)\).  Thus the maintained mutual-absolute-continuity member property is exactly \(Q\ll P\), proving
\[
 \mathfrak M_e(P)=\mathfrak M^{\to}_e(P)\cap\{Q:Q\ll P\}. \tag{14c}
\]
On a full-support binary alphabet, \(P\ll Q\) forces \(Q\) to have full support, hence also \(Q\ll P\).  Therefore the one-sided and mutual versions of each class coincide there.  Equations (14a)--(14c) prove the proposition.